
\section{Casi d'uso implementati}
In questa terza e ultima fase di sviluppo, l'obiettivo principale è stato il completamento della gestione del profilo utente e l'introduzione di funzionalità avanzate per l'esportazione e la visualizzazione dei viaggi in formato pdf. È stato inoltre effettuato un lavoro di rifinitura e miglioramento generale del frontend per garantire un'esperienza utente più fluida e intuitiva.

Di seguito sono riportati i casi d'uso sviluppati per questa iterazione, in conformità con i requisiti iniziali.

\begin{table}[htbp]
\centering
\renewcommand{\arraystretch}{1.3}
\begin{tabular}{|l|l|}
\hline
\textbf{ID} & \textbf{Titolo} \\ \hline
\textbf{UC6} & Esporta percorso \\ \hline
\textbf{UC2.1} & Modifica dati \\ \hline
\textbf{UC2.2} & Elimina account \\ \hline
\end{tabular}
\caption{Iterazione 3 - Casi d'uso implementati}
\label{tab:iterazione3_casi_uso}
\end{table}

\subsection*{UC6: Esporta percorso}
\textbf{Attori:} Utente, Sistema. \\
\textbf{Descrizione:} L'utente può esportare un itinerario in un formato condivisibile (PDF). Il sistema genera un documento contenente una tabella riassuntiva con l'ordine delle tappe, i nomi dei monumenti e i tempi di visita. Questa funzionalità attiva direttamente il modulo \textit{Export Service} previsto dall'architettura. \\
\textbf{Flusso degli eventi:}
\begin{enumerate}
    \item L'utente apre i dettagli di un viaggio salvato o appena ottimizzato.
    \item Clicca sull'opzione "Esporta in PDF".
    \item Il frontend invia la richiesta all'Application API.
    \item L'Export Service elabora i dati del viaggio e genera il file PDF impaginato.
    \item Il file viene scaricato automaticamente sul dispositivo dell'utente.
\end{enumerate}

\subsection*{UC2.1: Modifica dati}
\textbf{Attori:} Utente, Sistema. \\
\textbf{Descrizione:} L'utente autenticato può aggiornare le proprie informazioni personali e gestire la sicurezza del proprio account effettuando il cambio della password. \\
\textbf{Flusso degli eventi:}
\begin{enumerate}
    \item L'utente accede alla sezione dedicata alla gestione del profilo.
    \item Inserisce i nuovi dati personali o compila i campi per il cambio password (vecchia password e nuova password).
    \item Invia la richiesta di aggiornamento.
    \item Il sistema valida i dati immessi, aggiorna il database e restituisce un messaggio di conferma.
\end{enumerate}

\subsection*{UC2.2: Elimina account}
\textbf{Attori:} Utente, Sistema. \\
\textbf{Descrizione:} L'utente può richiedere la cancellazione definitiva e irreversibile del proprio profilo dal sistema. \\
\textbf{Flusso degli eventi:}
\begin{enumerate}
    \item L'utente accede alla sezione di gestione del profilo.
    \item Seleziona l'opzione "Elimina account".
    \item Il sistema richiede una conferma dell'operazione per evitare cancellazioni accidentali.
    \item L'utente conferma l'eliminazione.
    \item Il sistema rimuove in modo permanente i dati dell'utente dal database, chiude la sessione e reindirizza l'utente alla pagina di Login.
\end{enumerate}

